\begin{table}[htbp]\centering\scriptsize
\caption{\textbf{Table S5b.} Data partitions and sample accounting. Partitions are disjoint sets of annotated files (one annotation row per file); the unit separating them is the recording (file). Bats, channels and recording days overlap: all 33 positive-coded bats of the replication sample and of the completion sample, and all 12 channels, also occur in the discovery sample (recording days: 329 of 330 and 338 of 338 shared), so the out-of-sample sets test over-fitting to the discovery recordings, not generalisation to new individuals, recording set-ups or colonies.}\label{tab:S5b}
\begin{adjustbox}{max width=\linewidth}
\begin{tabular}{lrrrrrrp{5.2cm}}
\toprule
partition & files & Isolation files & positive-coded files & bats & channels & days & role \\
\midrule
discovery & 41,116 & 2,740 & 18,450 & 41 & 12 & 344 & all development, exploratory discovery of the cluster, supervised tests blocked by bat/channel/treatment/age, choice of settings \\
replication & 12,000 & 695 & 4,650 & 33 & 12 & 330 & out-of-sample test of C1--C3 (random draw) \\
completion & 33,557 & 2,089 & 13,344 & 33 & 12 & 338 & second out-of-sample test of C1--C3 (all remaining files) \\
reproduction only & 1,313 & 83 & 1,313 & 10 & 9 & 217 & extracted only to reproduce the June table; not used for any test \\
all annotated files & 87,986 & 5,607 & 37,757 & 41 & 12 & 345 & every annotated file of the corpus (205,252 further recordings are unannotated and unused) \\
\bottomrule
\end{tabular}
\end{adjustbox}
\end{table}
